\chapter{Introducción}

El melanoma cutáneo es la variante de cáncer de piel con mayor capacidad de diseminación metastásica, y su pronóstico clínico depende de su identificación en fases iniciales~\cite{amaral2025esmo}. Según las estimaciones mundiales de GLOBOCAN para 2024 (publicadas en julio de 2026), esta neoplasia registró 337.969 casos nuevos y 58.398 muertes en el mundo~\cite{sung2026globocan,iarc2026melanoma}. En Colombia, la Cuenta de Alto Costo informó que el 67,4\,\% de los casos diagnosticados en el periodo 2024 se identificó en estadios tempranos I y II, dentro de un reporte preliminar de 11.064 casos atendidos~\cite{cac2025melanoma}. Aunque la dermatoscopia realizada por especialistas mejora la precisión diagnóstica antes de la confirmación histopatológica obligatoria~\cite{amaral2025esmo}, la disponibilidad de herramientas tecnológicas de apoyo sigue siendo limitada fuera de entornos clínicos especializados~\cite{un2024sdg}.

El aprendizaje profundo ha cobrado relevancia en el análisis y clasificación de imágenes médicas~\cite{goodfellow2016deep,esteva2017dermatologist}. Las redes neuronales convolucionales (CNN) permiten extraer características morfológicas de imágenes dermatoscópicas y diferenciar entre lesiones benignas y malignas como apoyo a la valoración de los especialistas~\cite{esteva2017dermatologist,sariates2025}. Este tipo de clasificadores se ha evaluado incluso en estudios clínicos prospectivos como PROVE-AI (603 lesiones)~\cite{marchetti2023proveai} y en cohortes multicéntricas de 1.910 lesiones en ocho hospitales universitarios~\cite{heinlein2024multicenter}. De forma complementaria, el aprendizaje por transferencia (\textit{transfer learning}) adapta modelos preentrenados en ImageNet~\cite{deng2009imagenet} al dominio de las lesiones cutáneas, reduciendo los datos y el tiempo de entrenamiento requeridos~\cite{zhang2024dive,cherukuri2025,bello2024densenet}.

No obstante, persiste una brecha ingenieril y metodológica entre los modelos experimentales y su implementación en soluciones funcionales para el usuario~\cite{farias2026,almamun2025}. En el plano metodológico, el conjunto de referencia HAM10000 reúne 10.015 imágenes de 7.470 lesiones (con 1.956 lesiones fotografiadas más de una vez)~\cite{tschandl2018ham10000}, por lo que dividir los datos por imagen reparte fotografías de una misma lesión entre entrenamiento y prueba e infla las métricas por fuga de información~\cite{kapoor2023leakage,abhishek2025quality}. En el plano de despliegue, gran parte de las propuestas se limita a entornos de laboratorio y presenta barreras de usabilidad, privacidad o dependencia de servidores dedicados~\cite{farias2026,ning2024onnxweb}. El problema exige tanto entrenar bajo una partición libre de fuga como conectar el clasificador con una interfaz web desplegable, verificada y accesible~\cite{almamun2025,wang2025browser}.

El presente trabajo desarrolla un sistema de apoyo al diagnóstico de melanoma sobre imágenes dermatoscópicas que integra CNN y \textit{transfer learning} con un aplicativo web funcional desplegado en \pendiente{URL}. El sistema cubre el preprocesamiento de imágenes en cuatro etapas, el entrenamiento y evaluación de ocho configuraciones en 224 y 384\,px bajo una partición agrupada por lesión sobre HAM10000~\cite{tschandl2018ham10000}, la selección preregistrada de MobileNetV3Large~\cite{howard2019mobilenetv3} y su validación externa en el conjunto International Skin Imaging Collaboration (ISIC) 2018~\cite{codella2019isic2018}. Adicionalmente, incorpora la calibración de probabilidades~\cite{guo2017calibration} con modos de decisión configurables y su despliegue en una plataforma web que ejecuta el análisis en el navegador mediante OpenCV.js y ONNX Runtime Web~\cite{ning2024onnxweb}, donde los usuarios pueden cargar imágenes dermatoscópicas y obtener una predicción preliminar de apoyo sin enviarlas a servidores externos. El aporte del proyecto articula un protocolo de clasificación libre de fuga por lesión con una herramienta web accesible, conectando la experimentación en aprendizaje profundo con su despliegue práctico como apoyo a la valoración dermatoscópica~\cite{amaral2025esmo,farias2026}.

El resto del documento se estructura en ocho capítulos adicionales. El Capítulo~2 presenta la descripción del problema, la pregunta de investigación, la sistematización, los objetivos, la justificación vinculada al Objetivo de Desarrollo Sostenible (ODS) 3~\cite{un2024sdg} y las delimitaciones. El Capítulo~3 desarrolla el marco teórico sobre lesiones pigmentadas, dermatoscopia, CNN, \textit{transfer learning}, desbalance y métricas~\cite{maierhein2024metrics}, junto con los trabajos relacionados. El Capítulo~4 describe el conjunto HAM10000~\cite{tschandl2018ham10000}, la partición agrupada por lesión, el preprocesamiento, el aumento en línea y los pesos de clase; a continuación, el Capítulo~5 detalla las arquitecturas, la receta en dos fases, el experimento de fuga por unidad de partición, la regla de selección y la calibración~\cite{guo2017calibration}. Por su parte, el Capítulo~6 documenta la arquitectura del aplicativo web~\cite{ning2024onnxweb}, los requisitos, la interfaz, el manual de usuario y la verificación técnica~\cite{wang2025browser}. Finalmente, el Capítulo~7 expone los resultados en HAM10000 e ISIC 2018~\cite{codella2019isic2018}, la discusión y las limitaciones; el Capítulo~8 sintetiza las conclusiones por objetivo específico y el Capítulo~9 plantea el trabajo futuro~\cite{vasey2022decideai,daneshjou2022disparities,collins2024tripodai}.
